\chapter{Review of Related
Literature}\label{review-of-related-literature}

In the current landscape of blockchain technology, the emergence of
heterogeneous networks has revealed the term "information isolated
isalands" or "value silos" \citep{wang2021sok} \citep{Bigiotti2025}
\citep{Wu_2023}. Although individual blockchains proved to be useful as
independent ledgers, the lack of communication between them limits their
potential widespread adoption \citep{Hassan_2025} \citep{Cheng_2024}
\citep{Lohachab_2021}. Blockchain interoperability then becomes a
critical interaction between distributed systems \citep{wang2021sok}
\citep{Li2025} \citep{Lohachab_2021}.

Blockchain systems have been described as an independent state machine
with its own concensus rules and encryption standards
\citep{Bigiotti2025} \citep{Belchior_2021}. Because they do not have a
shared state or global clock between them, their interactions must
happen asynchronously; they must instead use message passing and state
synchronisation \citep{wang2021sok} \citep{Wei_2023}. In other words,
blockchain interoperability is a process wherein one source chain is
trying to change the state of a target chain by means of delivering
information or assets under certain securoty assumptions
\citep{wang2021sok} \citep{Wang_2023} \citep{belchior2021survey}.

Most of study in the literature focus on the blockchain interoperability
architecture. That is, whether it uses relays, sidechains, or notaries
\citep{Kotey_2023} \citep{Caetano_da_Silva_2025} \citep{Li_2023}
\citep{Chen_2023}. Nevertheless, recent studies have put emphasis that
architecture only gives stucture and real security and effectiveness of
interoperability depends on the behaviour of the system
\citep{Llambas2025} \citep{Llambas2024eventb} \citep{Lohachab_2021}.
This behaviour dictates how systems guarantee ACID (atomicity,
consistency, isolation, durability) properties in an asynchronous
network \citep{wang2021sok} \citep{Bigiotti2025} \citep{Wang_2023}
\citep{Li2025}.

This means that the sequencing of state transitions, timing of
cyrptographic proofs, and reactions to failure or timeout reveals the
system\textquotesingle s behaviour \citep{Ou_2022}
\citep{Chervinski_2023} \citep{deng2025} \citep{Wei_2023}.The analysis
of behaviours, instead of the static design, gives deeper understanding
on how different components (such as smart contracts and relayers)
interact during different processe. That is so to achieve semantic
interoperability where the meaning and context of data or asset from the
source chain gets accurately interpreted and valid to the target chain
\citep{Belchior_2023} \citep{sottomayor2024enterprise} \citep{Li2025}.

\section{Evolution of Blockchain
Interoperability}\label{evolution-of-blockchain-interoperability}

Surveying the literature, it is noticeable that the classifications of
interoperability reflects the underlying logic on how to trust and treat
information between blockchains \citep{wang2021sok} \citep{Pillai_2022}
\citep{wang2023}. Although categorical labels like notary schemes,
sidechains/relays, and hash-locking are usually used, there is a deeper
abstraction that surfaces in recent literature. That is, the perspective
of looking at blockchain interoperability as a multi-layered problem
derived from the OSI model of the internet \citep{Hassan_2025}
\citep{deng2025}.

The common abstraction from most of the sources is the
\textbf{source-target model}. In this model, an operation from the
source chain needs to trigger a corresponding state change to the target
chain \citep{deng2025} \citep{Cheng_2024}. The following differentiates
a critical difference on how this interaction gets seen:

\begin{itemize}
\item
  \textbf{Interoperability as Transfer:} In this perspective, the logic
  focuses on migration or transfer of ownership. Every transaction must
  follow all-or-nothing atomicity in order to avoid double-spending
  \citep{wang2021sok} \citep{Llambas2024patterns} \citep{deng2025}.
  Asset-based protocols belong here, ensuring that the value will not
  vanish or gets duplicated between different silos
  \citep{Belchior_2023} \citep{wang2023}.
\item
  \textbf{Interoperability as Communication:} In this perspective, the
  system gets treated as a \textbf{message-oriented service}
  \citep{wang2021sok} \citep{Kotey_2023} \citep{kang2022}. The goal is
  not to transfer the content of the source chain but to invoke a
  functionality in the target chain using arbitrary data
  \citep{Pillai_2022} \citep{Belchior_2021}. This is the logic of the
  data-based protocols like Arbitrary Message Passing (AMP)
  \citep{augusto2025}.
\end{itemize}

It is critical to make the distinction between asset-based and
data-based interoperability because they usually remain implicit in a
number of older literature. For instance, in the classical SoK
(Systematisation of Knowledge) of Buterin, the focus was always in value
exchange or value islands \citep{wang2021sok} \citep{A_Alzahrani_2025}
\citep{Pillai_2022}. However, in the newer frameworks, there is an
explicit division in the domain between \textbf{Asset Interoperability}
(migration/exchange) and \textbf{Data Interoperability} (cross-chain
calls/integration) \citep{deng2025} \citep{Cheng_2024} \citep{Li2025}.
This gap is important as a number of current models are designed for
assets but gets applied to data or vice-versa without taking each
different semantic requirements into consideration \citep{Belchior_2023}
\citep{Li2025}.

This differeng individual trust boundary and security assumptions of
each approaches contributes to the fragmented landscape of blockchain
interoperability\citep{Belchior_2023}. Although there is an attemot to
have a Generic Interoperability Paradigm comprising of Setup, Commit,
Verify, and Abort phases, the literature has not fully explained how
these phases differ in their \textbf{observable behaviour} when assets
are passed isnstead of arbitrary data \citep{deng2025} \citep{Li2025}.
This forms the need of a \textbf{comparative semantic analysis} to see
if these models are enough to encompass the behaviour of both asset and
data exchanges \citep{Li2025}.

\begin{itemize}
\tightlist
\item
  \textbf{Asset-Based Interoperability}
\end{itemize}

The literature treats asset-based interoperability not solely as a tool
to transfer currency but also as a \textbf{closed system of value
synchronisation} \citep{Pillai_2022} \citep{Cheng_2024}. The fundamental
logic of these protocols is the preservation of the INTEGRITY of assets
while they are moving between different ledgers
\citep{Llambas2024patterns} \citep{Pillai_2022} \citep{Belchior_2023}.

As a system, asset-based interoperability comprises the following
interdependent components: \textbf{Chains (Source and Target):} These
serve as state machines that record ownership \citep{wang2021sok}
\citep{deng2025} \citep{Cheng_2024}. \textbf{Smart Contracts (Custodian
vs. Debt Issuer):} In the source chain, there is usually a custodian
contract (vault) that locks asset; while in the target chain, there is a
debt issuer (money printer) that makes synthetic representation or
wrapped tokens \citep{Notland_2026} \citep{soureshjani2026}.
\textbf{Users and Intermediaries:} Included here is the initiator
(user), relayers that transfer proofs, and attesters that verify if the
lock event is valid before doing the minting \citep{Notland_2026}
\citep{Llambas2025} \citep{Belchior_2023}.

This shows that asset-based interoperability is a content-aware system.
Meaning, it is necessary that the protocol knows the specific semantics
of asset (e.g. ERC-20 standard) in order to ensure that the minted token
is semantically equal to the original asset \citep{Llambas2025}
\citep{Pillai_2022} \citep{deng2025}.

The observable behaviour of asset-based protocols follow a strict
sequential execution flow. Based on existing formal models (such as
Event-B analysis), the usual sequence follows: \textbf{Lock/Burn
(Source):} Immobilisation of the original value \citep{Notland_2026}
\citep{Llambas2025} \citep{deng2025}. \textbf{Verification:} Passing the
Merkle proof or SPV (Simplified Payment Verification) proceeding to the
target chain \citep{wang2021sok} \citep{lu2024atomicity}
\citep{Bigiotti2025} \citep{deng2025}. \textbf{Mint/Claim:} Producing of
an equal value at the target chain \citep{Llambas2025}
\citep{Sober_2022} \citep{Li2025}.

This shows that there is happens-before relationship. The mint
transaction at the target chain has causal dependency on the lock
transaction at the source chain \citep{Llambas2025} . If this sequence
is not followed, there is a risk of atomicity failure where there is a
possibility of minting assets that are not locked; resulting to
double-spending \citep{Llambas2024patterns} \citep{Belchior_2023}
\citep{deng2025} \citep{Li_2023}.

The major invariant that safeguards this system is the Global Supply
Invariant. This refers to maintaning a constant value to the overall
count of asset units in all networks \citep{Pillai_2022} \citep{Li_2023}
\citep{Bigiotti2025}.

It is critical that the safety (no bad behaviour or duplication) of
asset-based protocols gets prioritised over the liveness (fast
termination) \citep{Llambas2025} \citep{Bigiotti2025}. Because of this,
the observable behaviour has long confirmation periods (to wait the
finality from the source chain) and competition periods (for fraud
proofs) \citep{deng2025} \citep{wang2021sok} \citep{Wang_2023}
\citep{Mao_2023}. These delays are not just performance issue but a
behavioural necessity in order to protect the atomicity of the system
\citep{Chervinski_2023} \citep{lu2024atomicity} \citep{Cheng_2024}.

Nevertheless, this behaviour results to capital inefficiency
\citep{soureshjani2026} \citep{Mohammed_Abdul_2024} \citep{Li2025}. That
is because of the locking of assets, the liquidity gets trapped in the
contracts, that serves as a big limitation that is not usually seen in
data-based approaches \citep{soureshjani2026} \citep{deng2025}
\citep{augusto2025}. This creates the gap in the literature where even
though the safety of asset transfers can be formally modelled, there is
insufficient comparison if this strict ordering is also needed for
arbitrary data exchange \citep{Llambas2025} \citep{Llambas2024eventb}.

\begin{itemize}
\tightlist
\item
  \textbf{Data-Based Interoperability}
\end{itemize}

Unlike asset-based interoperability where it is illustrated as a
mechanism of changing ownership, the literature describes data-based
interoperability as a decentralised communication system
\citep{lu2024atomicity} \citep{Bigiotti2025} \citep{Kotey_2023}. If
asset-based focuses on liquidity, data-based on the other hand focuses
on the arbitrary message passimg (AMP). This is a mechanism where the
system acts as a transport layer for whatever type of information
\citep{Delgado_von_Eitzen_2025} \citep{augusto2025}.

Being a communication system, its architecture comprises of the
following interdependent components: \textbf{Chains and Modules:} Every
blockchain hosts modules or handlers (smart contracts). They serve as
entry and exit points of data \citep{Delgado_von_Eitzen_2025}
\citep{Chervinski_2023}. \textbf{Relayers and Connectors:} Since
blockchains are isolated islands and incapable of sending data natively,
they need off-chain relayers (such as Hermes) that scans blocks for
pending messages and transfer them to the target chain
\citep{Delgado_von_Eitzen_2025} \citep{Chervinski_2023}.
\textbf{Channels and Connections:} These are the stateful abstractions
that maintains logical link between two chains that gives ordering and
authentication to packets \citep{Bigiotti2025} \citep{wei2023formal}.

The literature sees the system as a payload-agnostic service. This shows
that the protocols do not care whether what is inside the packets is a
token transfer or a complex smart contract call. Its purpose is to
ensure that the opaque bytes are delivered with integrity
\citep{Kotey_2023} \citep{Delgado_von_Eitzen_2025}
\citep{Chervinski_2023} \citep{Bigiotti2025}.

The observable behaviour of data-based interoperability is guided by
handshake protocols and asynchronous execution \citep{lu2024atomicity}
\citep{Bigiotti2025} \citep{Wei_2023}. \textbf{Handshake Process:} This
is unlike the simple lock-and-mint; data-based systems follows a strict
4-step handshake (OpenInit -\textgreater{} OpenTry -\textgreater{}
OpenAck -\textgreater{} OpenConfirm) in order to establish a connection
\citep{Bigiotti2025} \citep{Wei_2023}. This behaviour is needed to
ensure that both chains are ready to accept data under the same security
assumptions \citep{Wei_2023}. \textbf{Asynchronous Message Passing:} The
interaction here is non-blocking. The source chain delivers a message
and can proceed to other operations while the packet is in-flight
\citep{lu2024atomicity} \citep{Bigiotti2025}. The success of the
interaction depends on Acknowledgements (ACK) or Timeouts. They are
observable events that dictates if an operation is deemed Delivered or
Expired \citep{Chervinski_2023} \citep{Wei_2023}.

The fundamental difference between the two is rooted on their
interaction logic: \textbf{State-Dependent vs. Instruction-Dependent:}
Unlike the asset-based that depends on state provenance (it needs the
verification of the lock state before minting), data-based behaviour is
instruction-dependent \citep{Wang_2022} \citep{deng2025}. Its aim is to
trigger a functionality in the target system that will work under its
own rules \citep{wang2021sok} \citep{Belchior_2021}. \textbf{Atomicity
vs. Eventual Delivery:} In asset transfers, atomicity is commonly
all-or-nothing under a capital constraint \citep{Pillai_2022}
\citep{Belchior_2023}. Meanwhile, in data exchange, it prioritises
liveness and delivery guarantees. It means that every message should
have a definite result (accepted or timeout) inside the asynchronous
network \citep{Bigiotti2025} \citep{Pillai_2022} \citep{Wei_2023}.

The behaviour here depends on the reliability of the transport layer.
Nonetheless, the critique in the current literature is the dependency of
these systems to untrusted layers \citep{Wei_2023}
\citep{Delgado_von_Eitzen_2025}. Formal analyses (such as the TLA+)
shows that even the standardized protocols can have trapped packets or
incomplete handshakes if the relayers\textquotesingle{} behaviour cannot
be properly modelled \citep{Wei_2023}. This pushed the need for deeper
comparative semantic analysis to see how it is possible to converge the
safety of assets and the liveness of data in one behavioural model.

\begin{longtable}[]{@{}lll@{}}
\caption{Survey and Protocol Design Literature}\tabularnewline
\toprule\noalign{}
Citation & Type of Work & Key Limitation \\
\midrule\noalign{}
\endfirsthead
\toprule\noalign{}
Citation & Type of Work & Key Limitation \\
\midrule\noalign{}
\endhead
\bottomrule\noalign{}
\endlastfoot
\citep{soureshjani2026} & Survey / Protocol Design (Asset) & No formal
model; no behavioural analysis \\
\citep{Notland_2026} & Survey / SoK (Hybrid bridges) & Lacks formal
system model; no concurrency semantics \\
\citep{Caetano_da_Silva_2025} & Survey (Cross-chain general) &
Architectural focus only; no behavioural equivalence \\
\citep{augusto2025} & Framework / Dataset (Data/AMP) & No
process-algebraic specification \\
\citep{Li2025} & Survey (Cross-chain general) & No formal comparison of
asset vs data behaviours \\
\citep{Delgado_von_Eitzen_2025} & Protocol Analysis (Data/IBC) & No
formal semantics; no LTS extraction \\
\citep{kishore2025interoperability} & Survey (Cross-chain general) &
Lacks behavioural criteria for unification/stacking \\
\citep{deng2025} & Survey (Asset + Data) & Identifies asset/data
distinction, no formal comparison \\
\citep{A_Alzahrani_2025} & Survey (General) & No concurrency or
bisimulation \\
\citep{Hassan_2025} & SoK (Layered architecture) & No process algebra;
static taxonomy only \\
\citep{Kumar_2025} & Protocol Design (Cross-chain transactions) & No
formal verification of behaviour \\
\citep{Bigiotti2025} & Survey (Data) & No behavioural semantics for
handshake protocols \\
\citep{Cheng_2024} & Protocol Design (Asset + Data) & Security
assumptions only; no behavioural model \\
\citep{sottomayor2024enterprise} & Framework (Semantic interoperability)
& No formal concurrency semantics \\
\citep{Mohammed_Abdul_2024} & Protocol Design (Cross-chain
communication) & No formal verification \\
\citep{Li_2024} & Survey (Data/logistics) & Domain-specific; no general
behavioural theory \\
\citep{Augusto_2023} & Protocol (Asset) & No formal semantics \\
\citep{Chen_2023} & Protocol Design (Asset/notary) & No behavioural
equivalence \\
\citep{Li_2023} & Survey (Cross-chain technology) & No formal
comparison \\
\citep{wang2023} & Systematic Survey (Asset + Data) & Identifies gap but
no formal solution \\
\citep{Belchior_2023} & Survey / Gap Analysis (Asset + Data) & Calls for
formal methods but does not apply PA \\
\citep{Wu_2023} & Protocol Design (Asset/notary + signatures) & No LTS
or bisimulation \\
\citep{Harris_2023} & Survey (Cross-chain challenges) & High-level
only \\
\citep{Kotey_2023} & Survey (Heterogeneous communication) & No formal
behaviour model \\
\citep{Mao_2023} & Survey (Cross-chain technology) & No concurrency
semantics \\
\citep{Wang_2022} & Protocol Design (Asset + Data) & No formal
verification \\
\citep{Sober_2022} & Protocol (Asset) & No behavioural equivalence \\
\citep{kang2022} & Survey (Interoperability landscape) & No formal
methods applied \\
\citep{Ou_2022} & Survey (Cross-chain mechanisms) & No process
algebra \\
\citep{Pillai_2022} & Framework (Asset/burn-to-claim) & No LTS or
branching behaviour \\
\citep{Lohachab_2021} & Survey (Interoperability role) & Behavioural
analysis missing \\
\citep{Belchior_2021} & Survey (Past, present, future trends) & No
formal semantics \\
\citep{wang2021sok} & SoK (Asset + Data) & Calls for behavioural
semantics but does not provide \\
\end{longtable}

\section{The IBC Protocol: A Unified Approach to
Interoperability}\label{the-ibc-protocol-a-unified-approach-to-interoperability}

While the previous sections have treated asset-based and data-based
interoperability as conceptually distinct categories, the
Inter-Blockchain Communication (IBC) protocol represents an attempt to
unify both within a single framework. Developed by the Cosmos ecosystem,
IBC has become the most widely adopted interoperability protocol,
connecting over 100 blockchains as of 2025 \citep{ibc2024transportapp}
\citep{Belchior_2021}.

\subsection{Protocol Architecture: Separation of Transport and
Application
Layer}\label{protocol-architecture-separation-of-transport-and-application-layer}

IBC is designed as a general-purpose message-passing protocol analogous
to TCP/IP for the Internet. Its architecture explicitly separates two
layers \citep{cosmos2024ibcintro} \citep{goes2020interblockchain}:

\begin{itemize}
\item
  \textbf{Transport, Authentication, and Ordering (TAO) Layer
  (ICS-004):} This layer handles the low-level concerns of packet
  delivery: establishing connections via a 4-step handshake (OpenInit →
  OpenTry → OpenAck → OpenConfirm), ordering packets, timeouts, and
  acknowledgements. The TAO is payload-agnostic -- it treats all data as
  opaque bytes \citep{Wei_2023} \citep{ibc2025howworks}.
\item
  \textbf{Application Layer (ICS-020, ICS-027, ICS-031):} This layer
  gives meaning to the payload. ICS-020 defines fungible token transfers
  (asset interoperability). ICS-027 defines Interchain Accounts
  (allowing one chain to control an account on another). ICS-031 defines
  Interchain Queries (reading state from another chain). The separation
  allows the same TAO to support both asset and data semantics
  \citep{quasar2023litepaper} \citep{ibc2024transportapp}.
\end{itemize}

\subsection{Asset Transfer in IBC
(ICS-020)}\label{asset-transfer-in-ibc-ics-020}

Under ICS-020, asset transfer follows a \textbf{lock-and-mint} pattern
consistent with asset-based interoperability described earlier:

\begin{enumerate}
\tightlist
\item
  \textbf{Lock/Burn on Source Chain:} The source chain locks the
  original tokens in a module account or burns them.
\item
  \textbf{Packet Transmission:} The lock event is packaged into an IBC
  packet and sent via the TAO layer.
\item
  \textbf{Mint on Target Chain:} Upon packet receipt and proof
  verification, the target chain mints an equivalent amount of vouchers
  (typically prefixed with the source chain\textquotesingle s channel
  identifier).
\item
  \textbf{Reverse Path (Transfer Back):} The same process in reverse
  burns vouchers and unlocks original tokens.
\end{enumerate}

This behaviour preserves the \textbf{Global Supply Invariant} or the
total supply across both chains remains constant \citep{Pillai_2022}
\citep{Li_2023}.

\subsection{Arbitrary Data Transfer in IBC
(ICS-004)}\label{arbitrary-data-transfer-in-ibc-ics-004}

At the TAO layer, the protocol exhibits data-based interoperability
semantics:

\begin{itemize}
\tightlist
\item
  \textbf{Asynchronous Packet Delivery:} Source chain commits a packet;
  target chain receives it independently. The source chain may continue
  other operations while the packet is in-flight
  \citep{lu2024atomicity}.
\item
  \textbf{Acknowledgement and Timeout:} Every packet must eventually
  result in either an acknowledgement (successful receipt and
  processing) or a timeout (failure to receive within a deadline). The
  alternative composition `ack + timeout` creates observable branching
  behaviour \citep{Chervinski_2023} \citep{Wei_2023}.
\item
  \textbf{Relayer-Mediated Forwarding:} Off-chain relayers scan both
  chains, detect packet commitments, and submit proofs to the target
  chain. Relayer behaviour is internal to the protocol -- from an
  external observer\textquotesingle s perspective, only packet send and
  receive/ack events are visible \citep{Delgado_von_Eitzen_2025}.
\end{itemize}

\subsection{Existing Formal Analysis of
IBC}\label{existing-formal-analysis-of-ibc}

To date, formal verification of IBC has focused primarily on the TAO
layer using trace-based methods: \textbf{Wei et al. (2023)} modelled
IBC\textquotesingle s TAO in TLA+ and discovered two critical
design-level issues: (1) \textbf{incomplete handshakes due to identifier
exhaustion} where channels could enter a state where no further
handshake steps were possible; (2) \textbf{trapped packets} where
packets could neither be received nor timed out under certain relayer
conditions \citep{Wei_2023}. While this analysis confirmed safety and
liveness properties of packet delivery, it left several questions
unanswered from a \textbf{behavioural semantics} perspective:

\begin{itemize}
\tightlist
\item
  Does the asset transfer behaviour (ICS-020) share the same core
  semantics as the data transport behaviour (ICS-004), or are assets
  simply stacked on top of data with additional constraints?
\item
  Do asset and data behaviours interfere with each other when both flow
  through the same channels?
\item
  What branching structures (e.g., timeout vs acknowledgement choices)
  characterise each behaviour, and are they equivalent?
\end{itemize}

These questions require a \textbf{behavioural} rather than trace-based
analysis. A behavioural analysis preserves the branching structure of
the system and can compare processes using bisimulation equivalence.

\subsection{The Need for Process-Algebraic Analysis of
IBC}\label{the-need-for-process-algebraic-analysis-of-ibc}

TLA+ verification, while valuable for detecting linear temporal logic
violations, models behaviour as sequences of states (traces). As will be
established in the theoretical framework, trace equivalence collapses
branching structure: `ack+timeout` (choice) is indistinguishable from
`ack.timeout` (sequence) under trace equivalence \citep{Fokkink2004}.

To answer whether IBC\textquotesingle s asset and data behaviours are
\textbf{unified}, \textbf{stacked}, or \textbf{conflicting}, a
process-algebraic approach is required. Process algebra can:

\begin{enumerate}
\tightlist
\item
  Model the concurrent composition of asset and data flows using the
  parallel operator
\item
  Preserve branching structure using alternative composition
\item
  Abstract internal relayer behaviour using the silent action
\item
  Compare behaviours using weak bisimulation rather than trace
  equivalence
\end{enumerate}

This research provides the first process algebra specification of
IBC\textquotesingle s ICS-004 and ICS-020 standards, filling the gap
identified above.

\section{Formal Methods to Blockchain
Interoperability}\label{formal-methods-to-blockchain-interoperability}

In the study of interoperability protocols, simple architecture diagrams
alone are insufficient to guarantee the security and correctness of
systems. Formal Methods surfaces in the literature as critical tools to
avoid design flaws that can lead to large financial losses
\citep{Wei_2023} \citep{Llambas2024eventb}. In surveying the literature,
formalisations of blockchain interoperability protocols have been
developed using a variety of methodologies such as interactive theorem
proving, model checking, and symbolic logic. They generally focus on
ensuring atomicity, safety, liveness, and consistency across
heterogeneous networks.

\begin{itemize}
\tightlist
\item
  \textbf{Mechanised Proofs and Refinement (Isabelle/HOL, Coq, Event-B})
  Isabelle/HOL has been used to provide rigorous, high-assurance proofs
  for several interoperability frameworks.

  \begin{itemize}
  \tightlist
  \item
    \textbf{Hierarchical Transactions (Canton)}: The Canton
    interoperability protocol\textquotesingle s transactions were
    formalised as authenticated data structures using a modular "Merkle
    functor" pattern \citep{lochbihler2020}. This approach allows for
    defining inclusion proofs and merging them into multi-proofs while
    maintaining the data integrity.
  \item
    \textbf{Global Consistency in IBC}: The IBC protocol was formalised
    as an instantiation of the Isabelle Infrastructure framework
    \citep{kammuller2020}. This model established a "global
    inconsistency" property, proving that scanning and submitting
    transactions between blockchains does not introduce inconsistencies
    across their respective ledgers.
  \item
    \textbf{Fail-Safe Cross-Chain Bridges}: A formal model of the Sonic
    Gateway bridge was verified to be safe using a stepwise refinement
    approach \citep{maric2025}. This verification proved that even if
    the receiver chain fails (a "dead bridge" scenario), users can still
    retrieve their assets on the sender chain via specific recovery
    transactions.
  \item
    \textbf{Cross-Domain Regulatory State Preservation}: A formal proof
    was established for regulatory state transitions (e.g., active or
    frozen status) across multiple domains \citep{kim2026}. It verified
    safety (bidirectional roundtrip preservation) and liveness
    (deterministic resolution of conflicting actions under Byzantine
    faults).
  \item
    \textbf{Bitcoin Validation Procedure}: The first formal model of
    Bitcoin\textquotesingle s transaction data structures and validation
    logic was mechanised in Coq \citep{rupic2020}. It provided
    machine-verified proofs for essential properties like the
    impossibility of double-spending and the uniqueness of transactions.
  \item
    \textbf{Inter-Contract Communication}: Coq has also been used to
    construct models that allow for interactions between contracts on
    different types of blockchains (e.g., depth-first vs. breadth-first
    execution models) \citep{kammuller2020}.
  \item
    \textbf{Gateway-Based Interoperability}: A gateway-based
    interoperability solution between Ethereum and Hyperledger Fabric
    was formalised in Event-B using the Rodin platform
    \citep{Llambas2024eventb}. It automatically discharged 23 proof
    obligations (17 invariant preservation, 6 well-definedness) using
    Rodin\textquotesingle s provers, confirming functional correctness
    and safety properties.
  \end{itemize}
\item
  \textbf{Temporal Logic of Actions (TLA+) and Model Checking} TLA+ and
  the TLC moderl checker have been utilised to identify functional and
  timing-related vulnerabilities in inter-chain protocols.

  \begin{itemize}
  \tightlist
  \item
    \textbf{IBC TAO Layer}: The TAO layer of IBC was modelled in TLA+
    \citep{Wei_2023} (see the preceeding section on IBC Protocol).
  \item
    \textbf{Tendermint Synchronisation (Fastsync)}: The Fastsync
    protocol was formalised at multiple levels (high-level, low-level,
    and concurrency) in TLA+ \citep{braithwaite2020}. Model checking
    helped refine properties related to termination and synchronisation,
    revealing that a global timeout might terminate the process before
    the node reaches the maximum height of its correct peers.
  \item
    \textbf{Cross-Chain Swaps}: Other efforts have included TLA+ formal
    proofs for specific cross-chain swap protocols and the
    "burn-to-claim" asset transfer protocol \citep{Wei_2023}.
  \end{itemize}
\item
  \textbf{Symbolic Logic and Tamarin}

  \begin{itemize}
  \tightlist
  \item
    \textbf{Hash Time Lock Contracts (HTLC)}: The security properties of
    HTLCs, the most common mechanism for atomic cross-chain swaps, were
    formally analysed using the Tamarin prover \citep{boyd2020}. This
    formalisation identified a critical "hidden" assumption: for the
    protocol to be secure, the growth speeds of the two blockchains
    involved must remain stable.
  \end{itemize}
\item
  \textbf{Timed Automata and UPPAAL}

  \begin{itemize}
  \tightlist
  \item
    \textbf{Pub-Sub Interoperability Protocol}: This protocol, which
    establishes communication between Hyperledger Fabric and Besu
    networks via a broker, was formalised as stochastic timed automata
    (STA) \citep{Alam_2023}. Using UPPAAL-SMC, researchers verified
    real-time properties and provided probabilistic estimations for
    message delivery, revealing an 8\% probability of failure when the
    broker\textquotesingle s throughput is significantly higher than
    that of the publishers and subscribers.
  \end{itemize}
\item
  \textbf{Abstract and Unified Frameworks}

  \begin{itemize}
  \tightlist
  \item
    \textbf{Atomic Multi-Chain Transaction Protocol}: Researches defined
    a uniform communication interface and a protocol inspired by
    two-phase commits to handle general transactions spanning several
    chains \citep{lu2024atomicity}. This was supported by formal
    modelling of blockchains as \textbf{abstract state machines} to
    prove atomicity.
  \item
    \textbf{SmartML and WhyISon}: Modelling languages like SmartML
    provide a platform-independent framework for analysing smart
    contracts \citep{veschetti2024smartmlmodelinglanguagesmart}.
    Similarly, WhyISon translates Michelson (Tezos) smart contracts into
    WhyML for automated deductive verification in the Why3 framework
    \citep{horta2020}.
  \item
    \textbf{QLink Architecture}: A quantom-safe bridge architecture was
    modelled mathematically to link cryptographic (Post-Quantom
    Cryptography) and physical-layer (Quantom Key Distribution) security
    \citep{silva2025}.
  \end{itemize}
\end{itemize}

Comparative Analysis: Process Algebra vs. Exisiting Formalisms

In contract to the state-based and trace-based approaches above,
\textbf{Process Algebras} (e.g., ACP, CSP, pi-calculus) offer a
compositionally rich behavioural perspective on interoperability. Below
is a summary of what each surveyed method lacks.

\begin{longtable}[]{@{}lll@{}}
\caption{Formal Methods Literature (Asset-Focused)}\tabularnewline
\toprule\noalign{}
Citation & Formal Method Used & Key Limitation \\
\midrule\noalign{}
\endfirsthead
\toprule\noalign{}
Citation & Formal Method Used & Key Limitation \\
\midrule\noalign{}
\endhead
\bottomrule\noalign{}
\endlastfoot
\citep{Llambas2025} & Event-B & State-based only; no concurrency or
branching \\
\citep{Llambas2024patterns} & Event-B & Lacks liveness and asynchronous
communication \\
\citep{Llambas2024eventb} & Event-B & No dynamic process creation; no
mobility \\
\citep{lochbihler2020} & Isabelle/HOL & No concurrency or branching
behaviour \\
\citep{maric2025} & Event-B refinement & State-based; no concurrency
semantics \\
\citep{rupic2020} & Coq & Single ledger; not interoperability-focused \\
\citep{boyd2020} & Tamarin & No true concurrency or liveness for data \\
\citep{silva2025} & Mathematical model & No process algebra; no
bisimulation \\
\citep{Wang_2024} & CSP (Process Algebra) & Single protocol; no
comparison of asset vs data \\
\end{longtable}

\begin{longtable}[]{@{}lll@{}}
\caption{Formal Methods Literature (Data-Focused and
General)}\tabularnewline
\toprule\noalign{}
Citation & Formal Method Used & Key Limitation \\
\midrule\noalign{}
\endfirsthead
\toprule\noalign{}
Citation & Formal Method Used & Key Limitation \\
\midrule\noalign{}
\endhead
\bottomrule\noalign{}
\endlastfoot
\citep{Wei_2023} & TLA+ & Trace-based only; no bisimulation or
branching \\
\citep{lu2024atomicity} & Abstract State Machines & No process algebra;
no bisimulation \\
\citep{kammuller2020} & Isabelle/HOL & No process algebra; trace-based
only \\
\citep{kim2026} & Isabelle/HOL & Safety/liveness but no branching
equivalence \\
\citep{braithwaite2020} & TLA+ & Not IBC-specific; no asset-data
comparison \\
\citep{Alam_2023} & Timed Automata (UPPAAL) & Probabilistic; no
behavioural refinement laws \\
\citep{horta2020} & Why3/WhyML & Single chain; not interoperability \\
\citep{Chervinski_2023} & Empirical / Measurement & No formal
specification \\
\end{longtable}

This gap motivates the need for a unified behavioural semantics --- one
that treats both asset and data interoperability as observable
interaction patterns rather than isolated state transitions.

\section{Synthesis and Research Gap}\label{synthesis-and-research-gap}

In surveying the literature, a clear pattern surfaces: although the
improvement in the technical solutions to blokchain interoperability is
growing fast, its theoretical and formal foundation still remains slow
and fragmented. The following lay outs the critical gaps that reasons
the need for this research.

\begin{enumerate}
\tightlist
\item
  \textbf{Gap of comparison}: No formal comparison of asset vs data
  behaviours within a single protocol.
\item
  \textbf{Gap of method}: No use of Process Algebra/LTS for
  interoperability semantics.
\item
  \textbf{Gap of behavioural criteria}: No operational definition of
  "unified/stacked/conflicting" in terms of observable behaviour.
\item
  \textbf{Gap of IBC semantics}: No process-algebra model of IBC that
  separates transport vs application behaviour.
\end{enumerate}

This research fills these gaps by providing the first process algebra
specificatio of IBC\textquotesingle s core interoperability standards
(ICS-004, ICS-020), extracting LTS models, and analysing whether asset
and data behaviours are unified, stacked, or conflicting.

\begin{longtable}[]{@{}lllll@{}}
\caption{Summary of Existing Work Against This
Research\textquotesingle s Key Criteria}\tabularnewline
\toprule\noalign{}
Reference & Formal Method & Asset-Data Comparison? & Uses Bisimulation?
& Models IBC? \\
\midrule\noalign{}
\endfirsthead
\toprule\noalign{}
Reference & Formal Method & Asset-Data Comparison? & Uses Bisimulation?
& Models IBC? \\
\midrule\noalign{}
\endhead
\bottomrule\noalign{}
\endlastfoot
\citep{wang2021sok} & None & No (calls for it) & No & No \\
\citep{Wei_2023} & TLA+ & No & No & \(\checkmark\) (TAO only) \\
\citep{Llambas2025} & Event-B & No & No & No \\
\citep{Llambas2024eventb} & Event-B & No & No & No \\
\citep{Llambas2024patterns} & Event-B & No & No & No \\
\citep{Wang_2024} & CSP & No (asset only) & Yes & No \\
\citep{lu2024atomicity} & ASM & No & No & No \\
\citep{kammuller2020} & Isabelle/HOL & No & No & Yes (consistency) \\
\citep{lochbihler2020} & Isabelle/HOL & No & No & No \\
\citep{maric2025} & Event-B & No & No & No \\
\citep{kim2026} & Isabelle/HOL & No & No & No \\
\citep{rupic2020} & Coq & No & No & No \\
\citep{boyd2020} & Tamarin & No & No & No \\
\citep{braithwaite2020} & TLA+ & No & No & No \\
\citep{Alam_2023} & Timed Automata & No & No & No \\
\citep{silva2025} & Mathematical & No & No & No \\
\citep{deng2025} & None & Yes (identifies gap) & No & No \\
\citep{Belchior_2023} & None & Yes (calls for formal) & No & No \\
\citep{Li2025} & None & Yes (calls for formal) & No & No \\
\citep{Hassan_2025} & None & No & No & No \\
\citep{Bigiotti2025} & None & No & No & No \\
\citep{Chervinski_2023} & Empirical & No & No & Yes (performance) \\
\citep{Delgado_von_Eitzen_2025} & Empirical & No & No & Yes
(empirical) \\
\citep{soureshjani2026} & None & No & No & No \\
\citep{Notland_2026} & None & No & No & No \\
\citep{Augusto_2023} & None & No & No & No \\
\citep{Pillai_2022} & None & No & No & No \\
\citep{Cheng_2024} & None & No & No & No \\
\citep{Kotey_2023} & None & No & No & No \\
\citep{Li_2023} & None & No & No & No \\
\citep{wang2023} & None & No & No & No \\
\citep{Wu_2023} & None & No & No & No \\
\citep{Chen_2023} & None & No & No & No \\
\citep{Sober_2022} & None & No & No & No \\
\citep{Wang_2022} & None & No & No & No \\
\citep{Mohammed_Abdul_2024} & None & No & No & No \\
\citep{Kumar_2025} & None & No & No & No \\
\citep{Caetano_da_Silva_2025} & None & No & No & No \\
\citep{augusto2025} & None & No & No & No \\
\citep{kishore2025interoperability} & None & No & No & No \\
\citep{A_Alzahrani_2025} & None & No & No & No \\
\citep{sottomayor2024enterprise} & None & No & No & No \\
\citep{Li_2024} & None & No & No & No \\
\citep{Harris_2023} & None & No & No & No \\
\citep{Mao_2023} & None & No & No & No \\
\citep{Ou_2022} & None & No & No & No \\
\citep{kang2022} & None & No & No & No \\
\citep{Lohachab_2021} & None & No & No & No \\
\citep{Belchior_2021} & None & No & No & No \\
\citep{horta2020} & Why3/WhyML & No & No & No \\
\textbf{\textbf{This research}} & \textbf{\textbf{ACP}} &
\textbf{\textbf{Yes}} & \textbf{\textbf{Yes}} & \textbf{\textbf{Yes}} \\
\end{longtable}
